\section*{Part IV: Economics, Size, and Strategy of the First Market}

\subsection*{44. Replace TAM Fascination with First-Stage Market Arithmetic}

If a founder plans to sell a course to working developers, a national estimate of education spending is nearly useless for planning the first cohort. The immediate question is how many qualified, reachable, able-to-pay people in the relevant situation exist during the actual enrollment window.

An illustrative bottom-up logic is:

\begin{equation}
\begin{aligned}
N_{\mathrm{addressable}} &= N_{\mathrm{eligible}} \times p_{\mathrm{trigger}} \times p_{\mathrm{reachable}}, \\
N_{\mathrm{buyers}} &\approx N_{\mathrm{addressable}} \times p_{\mathrm{consider}} \times p_{\mathrm{pay}}, \\
\mathrm{Gross\ bookings} &\approx N_{\mathrm{buyers}} \times \mathrm{realized\ fee}.
\end{aligned}
\end{equation}

This is a planning decomposition, not a validated behavioral law. Each factor must be grounded in a clearly defined period, avoid double-counting, and carry explicit uncertainty. Reachability and purchase probabilities are not necessarily independent. Some customers buy more than once; some cannot be identified in advance; some buy through partners. The expression is useful because it exposes hidden assumptions.

\subsection*{45. Worked Market Sizing for a Specialist Technical Studio}

Suppose, hypothetically, a founder identifies 1,200 potentially relevant engineers through ethical professional networks, communities, or public professional roles. The list is not a preauthorized mailing list; it is a research estimate, not permission to contact anyone. Assume 25\% currently face the specific triggering job, 40\% of that group can be reached in a legitimate relevant setting, 35\% seriously consider a defined offer, and 20\% of those pay.

\begin{table}[htbp]
\centering
\small
\begin{tabularx}{\columnwidth}{X r}
\toprule
\textbf{Funnel Stage} & \textbf{Hypothetical Count} \\
\midrule
Relevant population identified & 1,200 \\
Experiencing urgent trigger (25\%) & 300 \\
Reachable in buying window (40\%) & 120 \\
Seriously considering (35\%) & 42 \\
Paying customers (20\%) & 8.4 expected \\
\bottomrule
\end{tabularx}
\caption{Hypothetical bottom-up funnel for an engineering studio.}
\label{tab:ch7-studio-funnel}
\end{table}

At ₹35,000 of hypothetical realized fee, about 8 to 9 buyers would yield approximately ₹2.9 lakh of gross program sales, not profit. Small changes to the rates can halve or double the result. The founder should not call this a market forecast. It is a map of what to investigate and a way of evaluating whether the first cohort might be economically feasible.

If a program requires 20 paid students to break even and the realistic first reachable segment yields 8, the founder can change offer, price, timing, channel, fixed cost, or segment. Claiming a nationwide TAM cannot close the gap.

\subsection*{46. Bottom-Up Sizing for B2B Markets}

Now suppose an enterprise consultant targets a small class of technology firms. The questions differ:

\begin{itemize}
    \item How many eligible accounts can be named and plausibly reached?
    \item Which of them have the specific workflow or capability problem this year?
    \item Who controls the budget and when is the budget cycle?
    \item What fraction will engage an unfamiliar vendor?
    \item What is the likely scope and purchasing path?
    \item How much delivery capacity does 1 engagement absorb?
\end{itemize}

For example, 60 genuinely eligible accounts, 15 with an active qualifying project, 8 substantive discussions, 3 scoped pilots, and 1 paid agreement could be a useful entry market if that 1 contract is commercially sound and reproducible. But 1 large account can create founder dependence; the next deal may require a fresh 6-month cycle. Report both opportunity value and concentration risk.

\subsection*{47. Sizing Against the Business Model, Not the Founder's Ego}

A specialist may generate attractive income from a few customers. A venture-backed software company seeking hundreds of crores in revenue needs a much broader eventual demand system or very high per-customer value. A school seeking funding for research cannot assume a modest cohort surplus will directly purchase a campus.

Separate 3 questions:

\begin{enumerate}
    \item Can the first segment finance learning and survival?
    \item Can the initial operating system become repeatable and profitable?
    \item Can related segments or offerings compound enough value for the long-range ambition?
\end{enumerate}

The first does not establish the second or third. Conversely, the first market does not need to contain the entire long-term opportunity if the route into follow-on markets is explicit and plausible.

\subsection*{48. The Unit of Profitability Across Business Models}

Different industries have different economic atoms. For an online course, the economics may operate at the \textbf{cohort} level because scheduled teaching is shared while personal review costs vary with enrollment. For consulting, project contribution and partner utilization matter. For FMCG, 1 retailer can appear profitable until merchandising, spoilage, distributor fees, and repeat purchase are included. For marketplaces, an individual city and category may need enough dense transactions on both sides to become viable.

Your segment selection should use the unit on which cost and value actually accumulate. A startup cannot fix a broken per-cohort model by applying SaaS user-marginal-cost assumptions merely because the lessons are online.

\subsection*{49. A Worked Cohort Contribution and Capacity Constraint}

Assume illustratively a paid studio priced at ₹40,000 per learner, with ₹7,000 variable delivery and payment costs per participant, ₹4,000 acquisition costs per participant, and ₹1,80,000 fixed costs per cohort (including a normalized founder-teaching cost allocation). Then:

\begin{equation}
\Pi(n) = (₹40{,}000 - ₹7{,}000 - ₹4{,}000)n - ₹1{,}80{,}000.
\end{equation}

For 8 paying learners, the illustrative contribution after the fixed allocation is ₹52,000. For 12, it is ₹1,68,000. For 20, it is ₹4,00,000. But a cohort of 20 may exceed review capacity; that outcome is only possible if the advertised quality is preserved. At the same time, a founder with insufficient cash may not be able to risk the full fixed launch cost.

Track \textbf{cash contribution} and \textbf{normalized economic contribution} separately. Counting founder labor as free may be necessary for an initial cash experiment but can conceal an operating model that fails when a qualified instructor must be hired.

\subsection*{50. Margin Percentage Is Not Enough; Contribution Per Scarce Hour Matters}

A ₹3 lakh enterprise diagnostic may be attractive until it consumes 150 founder hours over 6 weeks. A ₹3 lakh cohort may require 55 founder hours plus paid mentors. Their cash margins and time economics differ.

One helpful measure is:

\begin{equation}
\begin{aligned}
&\text{Contribution per founder-hour} \\
&\quad= \frac{\text{Realized contribution}}{\text{Founder hours consumed}},
\end{aligned}
\end{equation}
interpreted alongside quality, skill development, and long-run value. A short engagement with poor margins can sometimes teach the company a unique repeated problem worth productizing. If that learning never produces reusable value, the engagement may simply be underpriced labor.

Founder time is not a literal cash expense when no salary is paid. It is a scarce resource with an opportunity cost. Recording it helps prevent a false claim of scalability.
